\section{Phạm vi đề tài}

Đồ án tập trung vào bài toán phát hiện trang bị bảo hộ lao động bằng YOLO26n
trên hai thiết bị đại diện là Jetson Nano và Jetson Orin Nano. Phạm vi tối ưu gồm
tìm siêu tham số, thưa hóa và cắt tỉa kênh, chưng cất tri thức, QAT, xuất ONNX
và tạo TensorRT engine. Các phép đo tập trung vào độ chính xác, latency, FPS, độ
phức tạp mô hình, bộ nhớ GPU và công suất.

Ở tầng nền tảng, đồ án xây dựng các chức năng cốt lõi: quản lý cluster, node,
runtime và deployment; điều khiển container tại biên; phân phối model; telemetry
và giám sát; Web Dashboard; AI Assistant. Hệ thống được triển khai thực tế với
Cloud trên một máy chủ AWS EC2 và các Edge Node thử nghiệm ở quy mô phòng thí
nghiệm (máy trạm x86 có GPU NVIDIA và thiết bị Jetson), chưa hướng đến thay thế
một nền tảng MLOps thương mại ở quy mô hàng nghìn node.

Đồ án không đặt mục tiêu nghiên cứu kiến trúc YOLO mới, phần cứng tăng tốc mới
hay mô hình ngôn ngữ nền tảng. Triển khai theo nhóm/canary, rollback tự động, mô
hình cảnh báo sớm bằng học máy và học liên tục được xác định là hướng phát triển.
